\section{Repositorio y equipo}
\label{sec:equipo}

\subsection{Repositorios}

\begin{table}[H]
\centering\small
\begin{tabularx}{\textwidth}{@{}l L@{}}
\toprule
Código de la app & \url{https://github.com/DesarrolloApps1/tpo-da1} \\
Documentación y diseño & \url{https://github.com/DesarrolloApps1/docs} (este documento y los archivos de Figma) \\
Backend & \pendiente{a crear} \\
\bottomrule
\end{tabularx}
\end{table}

\subsection{Forma de trabajo}

\begin{itemize}
\item \textbf{Ramas.} \texttt{main} tiene solo versiones entregables y está protegida. \texttt{dev} es la rama de integración. \texttt{tst} se usa para probar antes de pasar a \texttt{main}. Cada tarea se hace en \texttt{feature/<tema>}, que sale de \texttt{dev}.
\item \textbf{Commits.} Chicos y frecuentes, en español y con prefijo: \texttt{feat:}, \texttt{fix:}, \texttt{docs:}, \texttt{test:}, \texttt{refactor:}.
\item \textbf{Pull requests.} Un PR por funcionalidad hacia \texttt{dev}, con al menos una revisión de otro integrante y los tests pasando. Al cerrar un hito se pasa \texttt{dev} a \texttt{tst} y, después de probar, a \texttt{main}.
\item \textbf{Tareas.} Un issue por requisito o pantalla, organizados en un tablero de GitHub Projects.
\end{itemize}

\subsection{Integrantes, roles y responsabilidades}

Los roles no son exclusivos. Todos revisan código de otras áreas y tienen que poder explicar la app completa.

\begin{table}[H]
\centering\small
\begin{tabularx}{\textwidth}{@{}l l L L@{}}
\toprule
\textbf{Integrante} & \textbf{Rol principal} & \textbf{Responsabilidades} & \textbf{Participa en} \\
\midrule
\pendiente{Nombre} & Arquitectura e integración & Estructura de capas, navegación, AppContainer, revisión de PRs & Todas las áreas \\
\pendiente{Nombre} & UI y UX & Diseño en Figma, tema de Compose, pantallas y estados, accesibilidad & Pruebas con usuarios \\
\pendiente{Nombre} & Datos y sincronización & Room, DataStore, Retrofit, backend y estrategia offline & Canje \\
\pendiente{Nombre} & Detección y plataforma & Worker, UsageStatsManager, reglas de validación, notificaciones & Testing \\
\pendiente{Nombre} & Testing y documentación & Tests unitarios y de UI, README, este documento & Todas las áreas \\
\bottomrule
\end{tabularx}
\end{table}

\textbf{Cambios en el equipo:} \pendiente{ninguno hasta la fecha}.
